\documentclass[../../../include/open-logic-section]{subfiles}

\begin{document}

\olfileid{sth}{replacement}{extrinsic}
\olsection[اعتبارات خارجية]{اعتبارات خارجية بشأن الاستبدال}

نبدأ بالنظر في محاولة \emph{خارجية} لتسويغ الاستبدال. ويقترح بولوس
محاولة منها، على النحو الآتي:
\begin{quote}
  [\ldots] سبب تبني مسلّمات الاستبدال بسيط جدًا: فلها عواقب مرغوب
  فيها كثيرة، ولا عواقب غير مرغوب فيها (فيما يبدو). وإضافة إلى
  المبرهنات عن التصور التكراري، تشمل العواقب نظرية مرضية، وإن لم تكن
  مثالية، للأعداد اللانهائية، ونتيجة مرغوبًا فيها بشدة تسوّغ التعريفات
  الاستقرائية على العلاقات جيدة التأسيس.
  \citep[229]{Boolos1971}
\end{quote}
خلاصة فكرة بولوس أنه ينبغي لنا تسويغ الاستبدال بثماره. والثمار المعينة
التي يذكرها هي الأشياء التي ناقشناها في الفصول القليلة الماضية. فقد
أتاح لنا الاستبدال إثبات أن أعداد فون نويمان الترتيبية ممثلات ممتازة
لفكرة نمط الترتيب الجيد (وهذه هي «نظرية الأعداد اللانهائية المرضية،
وإن لم تكن مثالية» لدينا). وأتاح لنا الاستبدال أيضًا تعريف مجموعات
$V_\alpha$، وترسيخ مفهوم الرتبة، وإثبات الاستقراء بـ$\in$ (وهذا هو
مضمون «المبرهنات عن التصور التكراري» لدينا). وأخيرًا، يتيح لنا
الاستبدال إثبات مبرهنة العودية المتناهية التجاوز (وهذه هي «التعريفات
الاستقرائية على العلاقات جيدة التأسيس»).

هذه عواقب مرغوب فيها بالفعل. لكن هل تكفي هذه العواقب المرغوب فيها
لـ\emph{تسويغ} الاستبدال؟ \emph{لا}. أو لا تكفي، على الأقل، بصورة
مباشرة.

وفيما يأتي مشكلة بسيطة. فرغم أننا ذكرنا بعض عواقب الاستبدال المرغوب
فيها، كان يمكننا الحصول على كثير منها بوسائل أخرى. وهذا أمر ليس مشهورًا
بالقدر الذي ينبغي، ولذلك يلزم أن نتوقف لشرح الوضع.

ثمة نظرية بسيطة للمجموعات هي نظرية المستويات، أو اختصارًا~$\LT$.
\footnote{قدم~\citet{Montague1965} و\citet{Scott1974} أولى صور
  $\LT$؛ ثم بسّطها~\citet{Potter2004} وخصص لها معالجة في كتاب كامل؛
  وقد بسّط~\citet{ButtonLT1} مؤخرًا~$\LT$ أكثر.}
وليست مسلّمات~$\LT$ إلا الامتدادية والفصل والادعاء أن كل مجموعة
جزئية من \emph{مستوى} ما، حيث يعرّف «المستوى» بدهاء بحيث تسلك
المستويات مسلك أصدقائنا، مجموعات~$V_\alpha$. ومن ثم تثبت~$\ZF$
النظرية~$\LT$؛ لكن~$\LT$ \emph{أضعف بكثير} من~$\ZF$. والواقع أن
$\LT$ لا تعطينا الأزواج ولا مجموعات القوى ولا اللانهاية ولا الاستبدال.
ولتكن~$\Zr$ نتيجة إضافة اللانهاية ومجموعات القوى إلى~$\LT$؛ وهذا
يعطينا الأزواج أيضًا، فتكون~$\Zr$ قوية بقدر~$\Z$ على الأقل. لكن
$\Zr$ في الواقع أقوى تمامًا من~$\Z$، لأنها تضيف الادعاء أن لكل
مجموعة رتبة (ومن هنا اقتراحي أن نسميها~$\Zr$). والواقع أن~$\Zr$
تعطينا: نظرية مرضية تمامًا للأعداد الترتيبية؛ ونتائج تقسّم الهرم إلى
مراحل مرتبة ترتيبًا جيدًا؛ وبرهانًا على الاستقراء بـ$\in$؛
و\emph{صورة} من العودية المتناهية التجاوز.

وباختصار: مع أن بولوس لم يكن يعلم ذلك، كان يمكن الوصول إلى جميع
العواقب المرغوب فيها التي يذكرها \emph{من دون} الاستبدال؛ فلم يكن
يلزمه إلا استعمال~$\Zr$ بدلًا من~$\Z$.

(بالنظر إلى كل ذلك، لماذا اتبعنا الطريق التقليدي، فعلّمناك~$\ZF$
بدلًا من~$\LT$ و$\Zr$؟ ثمة سببان. الأول: لأسباب تاريخية محضة، ليس
البدء بـ$\LT$ معياريًا جدًا؛ وقد أردنا تزويدك بما يلزم لقراءة مناقشات
نظرية المجموعات الأكثر معيارية. والثاني: حين تصبح مستعدًا لفهم~$\LT$
و$\Zr$، تستطيع ببساطة قراءة~\citealt{Potter2004} و
\citealt{ButtonLT1}.)

وبالطبع، لما كانت~$\Zr$ أضعف تمامًا من~$\ZF$، فهناك نتائج تثبتها
$\ZF$ وتتركها~$\Zr$ مفتوحة. ولذلك قد يحاول المرء تسويغ الاستبدال
بأسباب خارجية بالإشارة إلى إحدى هذه النتائج. لكن من الصعب جدًا، متى
تعلمت كيفية استعمال~$\Zr$، العثور على أمثلة كثيرة لأشياء (أ) يحسمها
الاستبدال دون غيره، و(ب) تصدق حدسيًا. (للمزيد، انظر
\citealt[\S13.2]{Potter2004}.)

وخلاصة الأمر هي الآتية. لتقديم تسويغ خارجي مقنع للاستبدال، يلزم العثور
على نتيجة \emph{لا يمكن} تحصيلها من دونه. وليس ذلك بالمهمة السهلة.

لننظر في مشكلة أخرى تواجه كل محاولة لتقديم تسويغ خارجي محض للاستبدال.
(ولعل هذه المشكلة أعمق من الأولى.) لا يكتفي بولوس بالإشارة إلى أن
للاستبدال عواقب مرغوبًا فيها كثيرة؛ بل يقول أيضًا إنه لا يملك عواقب
غير مرغوب فيها «(فيما يبدو)». لكن هذا التحفظ بين قوسين، «فيما يبدو»،
حاسم قطعًا.

تذكر كيف بلغنا هذا الموضع: أدى الاستيعاب الساذج إلى عدم الاتساق،
فاستجبنا لذلك بتبني التصور التراكمي-التكراري للمجموعة. ويأتي هذا التصور
مصاحبًا لرواية نأمل أن تطمئننا إلى اتساقه. لكن إن لم نستطع تسويغ
الاستبدال من داخل تلك الرواية، فلا سبب لدينا (حتى الآن) للاعتقاد بأن
$\ZF$ متسقة. أو بتعبير أدق: لا سبب لدينا للاعتقاد بأن~$\ZF$ متسقة
غير الواقعة (العارضة ربما) أن أحدًا لم يكتشف تناقضًا \emph{بعد}. وبهذا
المعنى بالضبط يبدو تعليق بولوس راجعًا إلى هذا: «$\ZF$ متسقة (فيما
يبدو)». وينبغي لنا أن نطلب اطمئنانًا إلى الاتساق أكبر من ذلك.

ستؤثر هذه المسألة في كل محاولة \emph{خارجية محضة} لتسويغ الاستبدال،
أي كل تسويغ يصاغ فقط بعبارات عواقب~$\ZF$ (المعروفة). ولذلك نريد البحث
عن تسويغ \emph{داخلي} للاستبدال، أي تسويغ يوحي بأن الرواية التي
سردناها عن المجموعات تلزمنا «من قبل»، على نحو ما، بالاستبدال.

\end{document}
